\subsection{1. Ein gemeinsamer Träger, mehrere
Schatten}\label{ein-gemeinsamer-truxe4ger-mehrere-schatten}

Mit V=W†/√8 gilt V†V=I auf C\^{}60. Für eine logische Dichtematrix ρ
lautet der Schatten eines Fermionoperators O

\[\adjustbox{max width=\linewidth}{$\displaystyle 
  \Phi_{\mathcal O}(\rho)_O=\operatorname{tr}(\rho V^\dagger O V).
$}\]

Alle folgenden Messfamilien beziehen sich auf genau dieselbe ρ und
denselben Kodierer. Die Dimension der reellen Menge hermitescher
60×60-Matrizen ist 3600; bekannte Spur eins reduziert die freien
Zustandsparameter auf 3599. Alle angegebenen Familien enthalten die
Identität in ihrem linearen Raum, sodass ihre unsichtbaren Kerne bereits
spurfrei sind.

{\def\LTcaptype{none} % do not increment counter
\begin{longtable}[]{@{}
  >{\raggedright\arraybackslash}p{(\linewidth - 6\tabcolsep) * \real{0.2143}}
  >{\raggedleft\arraybackslash}p{(\linewidth - 6\tabcolsep) * \real{0.2857}}
  >{\raggedleft\arraybackslash}p{(\linewidth - 6\tabcolsep) * \real{0.2857}}
  >{\raggedright\arraybackslash}p{(\linewidth - 6\tabcolsep) * \real{0.2143}}@{}}
\toprule\noalign{}
\begin{minipage}[b]{\linewidth}\raggedright
Messfamilie
\end{minipage} & \begin{minipage}[b]{\linewidth}\raggedleft
Exakter Operatorrang
\end{minipage} & \begin{minipage}[b]{\linewidth}\raggedleft
Unsichtbare Zustandsdifferenzen
\end{minipage} & \begin{minipage}[b]{\linewidth}\raggedright
Verfügbarkeit
\end{minipage} \\
\midrule\noalign{}
\endhead
\bottomrule\noalign{}
\endlastfoot
Alle 64 Einzelbesetzungen n\_r & 24 & 3576 & Konkrete Operatoren;
ursprüngliche Messausführung offen \\
Alle hermiteschen Einteilchenbilineare f\_r†f\_s & 736 & 2864 &
Zusätzliche vollständige Bilinearfamilie \\
Alle G-gedrehten Einzelbesetzungen & 736 & 2864 & Bedingt auf
ausführbare G-Kontrollen \\
Alle 2016 ungedrehten Paarbesetzungen n\_i n\_j & 60 & 3540 & Gemeinsame
Besetzungsmessung erforderlich \\
Bilineare plus ungedrehte Paarbesetzungen & 772 & 2828 & Beide vorigen
Ressourcen zusammen \\
G-gedrehte Paarbesetzungen & 3600 & 0 & Gemeinsame Messung und
ausführbare G-Kontrollen \\
\end{longtable}
}

Der Rang zählt unabhängige lineare Erwartungswerte. Er zählt weder
Messapparaturen noch Messwiederholungen noch physische Raumrichtungen.

\subsection{2. Exakter Rangbeweis am
Tensor}\label{exakter-rangbeweis-am-tensor}

Die Datei enthält zehn symmetrische 16×16-Matrizen β\_k. Mit den sechs
antisymmetrischen 4×4-Farbmatrizen ε\_c gilt eintragsweise

\[\adjustbox{max width=\linewidth}{$\displaystyle 
 W_{(k,c),(p,a)(q,b)}=(\beta_k)_{pq}(\epsilon_c)_{ab}
 \quad\text{für }(p,a)<(q,b).
$}\]

Der Prüfer rekonstruiert hieraus den gesamten W und vergleicht alle
Einträge. Für sämtliche 4096 Bilineare vergleicht er zudem die direkte
CAR-Kontraktion mit der Faktorformel

\[\adjustbox{max width=\linewidth}{$\displaystyle 
 8V^\dagger f_{p,a}^\dagger f_{q,b}V
   =F_{pq}\otimes C_{ab},\qquad
 (F_{pq})_{kl}=\sum_t(\beta_k)_{pt}(\beta_l)_{qt},\quad
 (C_{ab})_{cd}=\sum_t(\epsilon_c)_{at}(\epsilon_d)_{bt}.
$}\]

In den aus den W-Faktoren abgeleiteten orthogonalen reellen Rahmen ist
der Spin-Bildraum genau \codeword{Skalar ⊕\allowbreak{} i·antisymmetrisch}, Dimension
1+45=46; der Farbbildraum hat Dimension 1+15=16. In der ursprünglichen
Gewichtsbasis prüft der Code die äquivalente ganzzahlige Identität

\[\adjustbox{max width=\linewidth}{$\displaystyle 
 n(A+\eta A^T\eta)=2\operatorname{tr}(A)I_n.
$}\]

Sie liefert jeweils die obere Ranggrenze. Ein nichtverschwindender
modularer Minor über dem Primkörper mit 1009 Elementen liefert die
passende untere Grenze. Damit sind die Ränge über Q und C exakt
bewiesen; es handelt sich nicht um eine toleranzabhängige numerische
Rangschätzung. Produktbildung ergibt 736.

Für die diagonalen Einzelbesetzungen bleiben nur Skalar plus Cartan:
(1+5)(1+3)=24. Die Adjungierten von so(10) und so(6) sind irreduzibel;
jede ursprüngliche Einzelbesetzung hat einen nichtverschwindenden
Skalar- und Cartananteil in beiden Faktoren. Ihr G-Orbit erzeugt deshalb
alle vier Summanden von (1⊕\allowbreak{}45)⊗(1⊕\allowbreak{}15), insgesamt 736, und keine
weiteren.

Der Eingabekern der gesamten Einteilchenkompression hat Dimension
\[4096-736=3360=(256-46)\cdot16.\] Er ist vom unsichtbaren \emph{Zustandskern}
mit Dimension 3600−736=2864 zu unterscheiden.

\subsubsection{Was genau fehlt?}\label{was-genau-fehlt}

Die gesamte Hermitesche Algebra zerfällt in den beiden reellen Faktoren
als

\[\adjustbox{max width=\linewidth}{$\displaystyle 
 (1\oplus i\Lambda^2\mathbb R^{10}\oplus\operatorname{Sym}^2_0\mathbb R^{10})
 \otimes
 (1\oplus i\Lambda^2\mathbb R^6\oplus\operatorname{Sym}^2_0\mathbb R^6).
$}\]

Die letzten symmetrischen, spurlosen Komponenten haben Dimension 54 und
20. Dem Bilinearschatten fehlen exakt die drei disjunkten Räume

\[\adjustbox{max width=\linewidth}{$\displaystyle 
 54\otimes(1\oplus15),\quad (1\oplus45)\otimes20,
 \quad54\otimes20,
$}\]

mit Dimension 864+920+1080=2864. Das ist eine konkrete Beschreibung des
vollständigen unsichtbaren Kerns, nicht bloß eine Dimensionsdifferenz.

Ein Verzicht auf alle imaginären Phasen wäre nochmals stärker: In
demselben reellen Gesamtrahmen haben reelle symmetrische Operatoren
Dimension 55·21+45·15=1830; imaginäre antisymmetrische Operatoren haben
Dimension 55·15+45·21=1770. Ein ausschließlich reeller Schatten könnte
diese 1770 Richtungen nicht sehen. Die native G-Familie enthält die
nötigen phasenempfindlichen Richtungen; ihre Ausführbarkeit wird damit
nicht bewiesen.

\subsection{3. Zwei orthogonale Zustände mit identischen eingeschränkten
Schatten}\label{zwei-orthogonale-zustuxe4nde-mit-identischen-eingeschruxe4nkten-schatten}

In der ursprünglichen logischen Gewichtsbasis, nullbasiert mit A=6k+c,
setze

\[\adjustbox{max width=\linewidth}{$\displaystyle 
 |\psi_\pm\rangle=\frac{|0\rangle\pm|30\rangle}{\sqrt2},
 \qquad \rho_\pm=|\psi_\pm\rangle\langle\psi_\pm|.
$}\]

Beide Matrizen sind exakt positiv, rein und zueinander orthogonal. Es
gilt für alle r,s und für alle i\textless j

\[\adjustbox{max width=\linewidth}{$\displaystyle 
 \operatorname{tr}[(\rho_+-\rho_-)V^\dagger f_r^\dagger f_sV]=0,
 \qquad
 \operatorname{tr}[(\rho_+-\rho_-)V^\dagger n_in_jV]=0.
$}\]

Insbesondere ändern beliebig viele G-gedrehte Einteilchenmessungen daran
nichts. Die Zustände unterscheiden sich in einer fehlenden
Spin-54-Richtung.

Der normierte logische Vektor

\[\adjustbox{max width=\linewidth}{$\displaystyle 
 |z\rangle=\tfrac12(|0\rangle+|30\rangle+i|6\rangle+i|36\rangle)
$}\]

liegt hingegen im G-Orbit eines ursprünglichen Kanalvektors. Die
zugehörige gedrehte Paarbesetzung komprimiert zu \textbar z⟩⟨z\textbar/8
und liefert die exakt verschiedenen Erwartungswerte 1/16 und 0. Dies ist
ein positiver Zustandszeuge für die Verbesserung durch die zusätzliche
Paarressource.

\subsection{4. Endliche, minimale Familie für vollständige
Rekonstruktion}\label{endliche-minimale-familie-fuxfcr-vollstuxe4ndige-rekonstruktion}

Jede unterstützte Paarspalte des W hat genau einen Eintrag±1. Daher gilt

\[\adjustbox{max width=\linewidth}{$\displaystyle 
 V^\dagger n_i n_jV=\tfrac18|A\rangle\langle A|
$}\]

für den zugehörigen Kanal A. Jeder der 60 Kanäle tritt auf. Ungedrehte
Paarbesetzungen erzeugen also exakt alle Diagonalmatrizen. Ihr Schnitt
mit dem Bilinearraum hat Dimension 24: Diagonalprojektion erhält die
oben angegebenen Faktor-Bildräume und hat in ihnen Rang 6 und 4. Damit
folgt der kombinierte Rang 736+60−24=772.

Der Prüfer leitet die reellen Rahmen aus den invarianten Faktormetriken
ab. Er prüft außerdem das W-Intertwining für sämtliche 45+15
Lie-Generatoren und ihre vollständigen Bilder so(10), so(6). G wirkt im
logischen Raum folglich als SO(10)×SO(6), jeweils mit seiner
Spin/SU(4)-Überlagerung.

Ein ursprünglicher Gewichtskanal ist in jedem Faktor ein zirkularer
Vektor (e\_i+i e\_j)/√2. Sein SO(n)-Orbit besteht aus (u+i v)/√2 mit
reellen, orthonormalen u,v. Für n=10 und n=6 erzeugt der Prüfer eine
rationale Liste solcher Projektoren:

\begin{enumerate}
\def\labelenumi{\arabic{enumi}.}
\tightlist
\item
  Alle (e\_i±i e\_j)/√2.
\item
  Für jedes i\textless j und ein drittes k den Vektor
  \codeword{(3e\_i+4e\_j+5i e\_k)/sqrt(50)}.
\end{enumerate}

Alle 50P haben ganzzahlige reelle und imaginäre Einträge. Der Code prüft
\codeword{P²=P}, \codeword{tr(P)=1} und die Isotropie
\codeword{tr(P Pᵀ)=0} exakt nach Skalierung. Die Isotropie und die Norm
kennzeichnen diese SO(n)-Gewichtsorbits; durch Orientierungswahl auf dem
orthogonalen Komplement existiert eine Drehung mit Determinante+1.
Beliebige U(10)- oder U(6)-Kontrollen werden nicht benutzt.

Modulare Zeilenauswahl liefert folgende konkrete unabhängige Basen:

\[\adjustbox{max width=\linewidth}{$\displaystyle 
 \mathcal B_{10}=\{I_{10},P_1,\ldots,P_{99}\},\qquad
 \mathcal B_6=\{I_6,Q_1,\ldots,Q_{35}\}.
$}\]

Die ausgewählten Vektoren und Indizes stehen vollständig im
Ergebnis-JSON. Ihre Produkte bilden eine Basis aller hermiteschen
60×60-Matrizen. Bei bekannter Spur bleibt eine minimale Familie von 3599
nichtkonstanten Zahlen:

\begin{itemize}
\tightlist
\item
  99·35=3465 gedrehte Paarprojektoren P\_a⊗Q\_b;
\item
  99 Spinwerte P\_a⊗I\_6, als Summe über die sechs Farbkanäle;
\item
  35 Farbwerte I\_10⊗Q\_b, als Summe über die zehn Spinkanäle.
\end{itemize}

Alle physisch angesetzten Paarerwartungswerte besitzen denselben Faktor
1/8; die Rekonstruktion kalibriert ihn ausdrücklich zurück. Da ein
offener Bereich um I/60 alle 3599 spurlosen hermiteschen Richtungen
enthält, kann keine feste Familie linearer Skalarerwartungswerte mit
weniger als 3599 Zahlen jeden Zustand identifizieren.
\textbf{Minimalität bezieht sich genau auf diese Zahlenzahl}, nicht auf
Geräteeinstellungen oder Messkosten.

\subsubsection{Unabhängige
Rekonstruktionsprüfung}\label{unabhuxe4ngige-rekonstruktionspruxfcfung}

Die Familie wird ohne Zustandsdaten festgelegt. Anschließend werden die
beiden obigen reinen Zustände sowie sechs deterministisch zufällige
Zustände der Ränge 1, 2, 5, 17, 60, 60 in derselben ursprünglichen
W-Codebasis verwendet. Aus ihren 3599 kalibrierten Schatten
rekonstruiert der Prüfer die vollständigen Dichtematrizen. Zusätzlich
sagt er je zwölf vorher unbenutzte, G-gedrehte Paarerwartungswerte
dieser selben Zustände voraus.

Diese Tests sind \textbf{numerisch}, nicht als exakte Zustandsbeweise
ausgegeben. Im dokumentierten Lauf beträgt der größte
Matrixeintragsfehler unter 7.6×10⁻¹⁶; der größte Fehler der 96
unbenutzten Paarvorhersagen liegt unter 5.8×10⁻¹⁷. Die Antwortmatrizen
der Faktoren haben Konditionszahlen ungefähr 40.47 und 20.72. Das ist
ein rauschfreier Rekonstruktionstest; statistische Messkosten,
Zustandspräparation und Robustheit gegen reales Rauschen sind hier nicht
nachgewiesen.

\subsection{5. Was native Zeitentwicklung
hinzufügt}\label{was-native-zeitentwicklung-hinzufuxfcgt}

Der gemeinsame helle N=2-Prozess besitzt H=h⊗I\_60. Für jeden
Zweigoperator B und inneren Operator O gilt exakt

\[\adjustbox{max width=\linewidth}{$\displaystyle 
 [h\otimes I,B\otimes O]=[h,B]\otimes O.
$}\]

Beliebige Zeitableitungen und Zeiten verändern daher nur den
Zweigfaktor. Sie erweitern den inneren Schattenrang 24 beziehungsweise
736 nicht. Die zusätzlichen Paarsonden können den inneren Rang dagegen
bis 3600 erhöhen. Zweigkohärenz und ihre Rekonstruktion erfordern eine
gesonderte Prozessbetrachtung; eine innere Zustandsrekonstruktion allein
bestimmt noch keine vollständige Mehrzeitgeschichte oder ein
Messinstrument.

\subsection{Reproduktion und Grenze}\label{reproduktion-und-grenze-1}

Mit Python 3.10 oder neuer und NumPy, aus diesem Ordner:

\begin{verbatim}
python3 -B replay.py
\end{verbatim}

Der Lauf überprüft die fünf eingefrorenen Quellenpins, führt den eigenen
Prüfer normal und mit \codeword{-OO} aus und verlangt byte-identische
Ausgaben. Warnungen gelten als Fehler; keine Prüfbedingung benutzt
abschaltbare Python-Assertions. Die Prüfzahlen stehen im
Replay-Manifest. Die vier Kontextdateien werden gehasht, nicht als
fremde Gesamtsuite erneut ausgeführt. Die mathematische Prüfung benutzt
nur die lokale W-Archivdatei.

Damit ist eine gemeinsame, am Ursprungstensor kalibrierte \emph{innere}
Schattenkarte konstruiert. Ein räumlicher Rand, ursprüngliche
ausführbare Instrumente, eine universelle Prozessidentifikation,
3+1D-Raumzeit oder ein Abschluss von T1-T8 folgt daraus nicht.

\begin{center}\rule{0.5\linewidth}{0.5pt}\end{center}

\section{Teilbericht C: kausaler
Eingriff}\label{teilbericht-c-kausaler-eingriff}

\section{Exakter kausaler Eingriff am unveränderten nativen
Hamiltonoperator}\label{exakter-kausaler-eingriff-am-unveruxe4nderten-nativen-hamiltonoperator}

\subsection{Ergebnis}\label{ergebnis-3}

Im ursprünglichen TFPT-Modell gibt es einen vollständig gelösten
kausalen Test innerhalb des Ladungssektors N=2. Ein Phasenimpuls an
Fermionmode 4 ändert die spätere unbedingte Besetzung von Mode 5. Die
beiden Observablen kommutieren vor der Entwicklung. Der Impuls erhält
die Ladung und wirkt nur auf Mode 4; es gibt keine Konditionierung auf
Messergebnisse.

Der stärkste hier bewiesene Vergleich endet in \textbf{beiden} Armen mit
genau zwei Fermionen und keinem Boson. Gemessen wird eine andere
Verteilung der Fermionen auf disjunkte Paare. Die natürliche Entwicklung
zwischen Eingriff und Messung verwendet ausschließlich das ursprüngliche
H und W.

\textbf{Bedingung:} Die Präparation des angegebenen Produktzustands, ein
gezielter Phasenimpuls und die einzelne Besetzungsmessung werden als
Instrumente gewährt. Ihre physische Verfügbarkeit aus dem ursprünglichen
Operationssatz ist hier nicht hergeleitet. Der Impuls ist eine
ausdrücklich hinzugewährte Kontrolloperation; er wird nicht als von der
ungestörten H-Entwicklung bereitgestellt ausgegeben. Es wurde kein
zusätzlicher Hopping- oder Linkterm in die freie Dynamik eingeführt.

Das Ergebnis betrifft Moden innerhalb derselben Bank. Es enthält weder
eine räumliche Metrik noch den bereits untersuchten N=64-Grundzustand
oder dessen geladenen Entnahmepol.

\subsection{1. Quelle und vollständiger invarianter
Sektor}\label{quelle-und-vollstuxe4ndiger-invarianter-sektor}

Die eingefrorene Quelle ist \codeword{sources/spinor\_tensors.npz} mit
SHA-256
\codeword{3f00a0892157a691e4ef26920c3702772c7bf76a250fad98ff04a079bb64b763}.

Es gilt unverändert

\[\adjustbox{max width=\linewidth}{$\displaystyle 
H=\Delta N_b+g\sum_A(b_A^\dagger P_A+P_A^\dagger b_A),\qquad
P_A=\sum_{i<j}W_{A,ij}f_jf_i,\qquad N=N_f+2N_b.
$}\]

W besitzt 60 Zeilen, 2016 Spalten und 480 Einträge mit Wert ±1. Jede
aktive Spalte gehört genau einer Zeile. Jede Zeile enthält acht
Fermionenpaare mit insgesamt sechzehn verschiedenen Moden. Deshalb
zerfällt der \textbf{gesamte} N=2-Sektor in 60 invariante
neundimensionale Blöcke und 1536 vollständig dunkle Paarzustände. Die
sieben dunklen Richtungen innerhalb jedes Blocks kommen zusätzlich
hinzu: insgesamt bleibt die bekannte dunkle Dimension 1536+60·7=1956.

Eine solche Gruppe besteht aus einem Boson und acht Paaren. Mit
\(|p_\ell\rangle=f_{i_\ell}^\dagger f_{j_\ell}^\dagger|0\rangle\),
\(s_\ell=W_{A,i_\ell j_\ell}\) und
\(|S\rangle=\sum_\ell s_\ell|p_\ell\rangle/\sqrt8\) gilt

\[\adjustbox{max width=\linewidth}{$\displaystyle 
H_A=\begin{pmatrix}0_{8\times8}&g s\\g s^T&\Delta\end{pmatrix},
\qquad H_A|_{\operatorname{span}(S,b_A)}=
\begin{pmatrix}0&\sqrt8g\\\sqrt8g&\Delta\end{pmatrix}.
$}\]

Die CAR-Vorzeichen wurden für alle 480 aktiven Paare separat geprüft.
Dies ist eine exakte Einschränkung auf einen invarianten Sektor, keine
Abschneidung höherer Bosonstufen.

Die benutzte Zeile A=0 lautet:

{\def\LTcaptype{none} % do not increment counter
\begin{longtable}[]{@{}rlr@{}}
\toprule\noalign{}
Blatt & Fermionenpaar & Vorzeichen \\
\midrule\noalign{}
\endhead
\bottomrule\noalign{}
\endlastfoot
0 & (4,57) & −1 \\
1 & (5,56) & +1 \\
2 & (8,53) & +1 \\
3 & (9,52) & −1 \\
4 & (16,45) & −1 \\
5 & (17,44) & +1 \\
6 & (28,33) & +1 \\
7 & (29,32) & −1 \\
\end{longtable}
}

